% ============================================================
% 分片：s02d  对应原文 PDF 第 15--17 页
% 负责章节：§2.2.3 Landau--Ginzburg models
% 译者：translator-03   审校状态：初稿
% ============================================================

% ==== 原文 PDF 第 15 页 ====
\subsubsection{Landau--Ginzburg 模型}

超空间除了能简洁地表示超多重态之外，还是把明显超对称的作用量泛函写成超空间积分的有用工具。例如，若 $K(S(z,\bar z;\theta^\pm,\bar\theta^\pm))$ 是超场 $S$ 的任意可微函数，则有
\begin{equation}
\int d^2z\,d^4\theta\,K\big(S(z,\bar z;\theta^\pm,\bar\theta^\pm)\big)\,,
\tag{2.21}
\end{equation}
其中 $\int d^4\theta\,K$ 取出 $K$ 的 $\theta^+\theta^-\bar\theta^+\bar\theta^-$ 项，此式自动超对称不变：$K$ 在 $Q_\pm$（相应地 $\bar Q_\pm$）变分下，第一项消去一个 $\theta^\pm$（相应地 $\bar\theta^\pm$），于是在费米子坐标的积分 $\int d^4\theta$ 下为零；第二项在费米子积分后存活，但结果是全导数，于是在对玻色子坐标积分后为零。这类表达式称为 $D$ 项。

% [Tutor Note: (2.21) 的"自动不变"与 $K$ 的具体形式无关，关键是超空间积分测度的性质：$Q_\alpha$ 在超空间中至多差一个全 $z$ 导数地等于一个"超空间导数"，故 $\delta_\varepsilon K$ 对 $\theta$ 是导数，$\int d^4\theta\,\delta K=0$；而 $\bar Q_\pm$ 变分留下的项在费米子积分后只剩 $\partial_z$ 全导数，对 $\int d^2z$ 亦为零。这也解释了为什么二维 $\mathcal{N}=(2,2)$ 理论可以只靠 $D$ 项与 $F$ 项两类超空间积分写出全部超对称作用量。]

% [Tutor Note: 物理图像：$D$ 项是二维 $\mathcal{N}=(2,2)$ 手征多重态的动能项，其几何意义是靶空间 $M$ 上非线性 $\sigma$ 模型的度规拉回；而 $F$ 项给出全纯超势，即靶空间上的全纯函数。二者合起来正是 (2.28) 所示的二维 $\mathcal{N}=(2,2)$ 理论的一般作用量，也是 $A$ 模型（只依赖凯勒结构）与 $B$ 模型（只依赖复结构）能分别从同一理论中扭变出来的原因。]

值得注意的是，若 $D$ 项只是手征多重态（或只是反手征多重态）的函数，则所得 $D$ 项为全导数。更一般地，在由（玻色型）手征多重态 $\Phi^n$（$n=1,\dots,N$）及其共轭反手征多重态 $\bar\Phi^{\bar n}$ 参数化某个复流形 $M$ 的理论中，$K$ 可解释为凯勒势。把凯勒势作如下移动，即加上手征场 $\Phi^n$ 的全纯函数及其共轭：
\begin{equation}
K(\Phi^n,\bar\Phi^{\bar n})\to K(\Phi^n,\bar\Phi^{\bar n})+\delta(\Phi^n)+\bar\delta(\bar\Phi^{\bar n})
\tag{2.22}
\end{equation}
这是一次凯勒变换：二维 $\mathcal{N}=(2,2)$ 手征多重态的复靶空间自然就是一个凯勒流形！显式完成费米子积分后，这样的 $D$ 项给出
% [Tutor Note: 凯勒变换 (2.22) 的关键在于 $\delta(\Phi^n)$ 只依赖手征场、$\bar\delta(\bar\Phi^{\bar n})$ 只依赖反手征场，因此二者之和不改变凯勒度规 $g_{n\bar n}=\partial_{\phi^n}\partial_{\bar\phi^{\bar n}}K$。物理上 $K$ 不是靶流形 $M$ 上的函数，而只是 $M$ 上某线丛的局部截面，允许相差一次凯勒变换；真正的几何数据是度规 $g_{n\bar n}$，这正是 $D$ 项能给出靶空间 $M$ 上非线性 $\sigma$ 模型的原因。]
\begin{equation}\tag{2.23}
\begin{aligned}
\int d^2z\,d^4\theta\, K(\Phi^n,\bar\Phi^{\bar n}) = \int d^2z\bigg[& g_{n\bar n}\big(\partial_{\bar z}\phi^n\partial_z\bar\phi^{\bar n}+\partial_z\phi^n\partial_{\bar z}\bar\phi^{\bar n}\\
&\quad +i\bar\psi^{\bar n}_-\mathcal{D}_{\bar z}\psi^n_- - i\bar\psi^{\bar n}_+\mathcal{D}_z\psi^n_+\big)\\
&+g_{n\bar n}\big(F^n-\Gamma^n{}_{ml}\psi^m_+\psi^l_-\big)\big(\bar F^{\bar n}-\Gamma^{\bar n}{}_{\bar m\bar l}\bar\psi^{\bar m}_-\bar\psi^{\bar l}_+\big)\\
&\quad +R_{n\bar n m\bar m}\psi^n_+\psi^m_-\bar\psi^{\bar n}_-\bar\psi^{\bar m}_+\bigg],
\end{aligned}
\end{equation}
其中 $g_{n\bar n}:=\partial_{\phi^n}\partial_{\bar\phi^{\bar n}}K$ 是凯勒度规；$\Gamma^n{}_{ml}$ 与 $\Gamma^{\bar n}{}_{\bar m\bar l}$ 是 $g_{n\bar n}$ 的 Christoffel 符号；$\mathcal{D}_z$、$\mathcal{D}_{\bar z}$ 是相应 Levi-Civita 联络的拉回，例如
% [Tutor Note: (2.23) 的第一行是动能项：$g_{n\bar n}$ 是靶空间度规的拉回，$\partial\phi\,\partial\bar\phi$ 是玻色子的动能，$i\bar\psi\mathcal{D}\psi$ 是费米子的协变动能（$\mathcal{D}_z$ 为 Levi-Civita 联络沿映射 $\phi$ 的拉回）。第二行含辅助场的双线性 $g_{n\bar n}(F^n-\Gamma\psi\psi)(\bar F^{\bar n}-\Gamma\bar\psi\bar\psi)$ 与四费米子项 $R_{n\bar n m\bar m}\psi\psi\bar\psi\bar\psi$。后者不是可选项：辅助场经运动方程消去后，标量势与四费米子相互作用的系数必须由靶空间 Riemann 张量给出，否则超对称性被破坏。]
\begin{equation}\tag{2.24}
\mathcal{D}_z\psi^n_+=\partial_z\psi^n_++\Gamma^n{}_{ml}(\partial_z\phi^m)\psi^l_+\,,\qquad
\mathcal{D}_{\bar z}\psi^n_-=\partial_{\bar z}\psi^n_-+\Gamma^n{}_{ml}(\partial_{\bar z}\phi^m)\psi^l_-\,;
\end{equation}
而 $R_{n\bar n m\bar m}$ 是 $g_{n\bar n}$ 的 Riemann 张量。

若把 $\phi,\bar\phi$ 解释为世界面 $\Sigma$ 到凯勒靶空间 $M$ 的映射 $\phi:\Sigma\to M$，其余场也都有简洁的几何描述。\footnote{以下描述源自如下设定：手征超场 $\Phi^n$ 按复靶空间上的全纯坐标变换。特别地，在坐标变换 $\phi^n\to\phi'^n(\phi)$ 下，手征超场（在移位后的坐标中）如下变换：
\[
\begin{aligned}
\Phi'^n(\Phi)&=\phi'^n(\phi)+\theta^+\Big(\frac{\partial\phi'^n}{\partial\phi^m}\,\psi^m_+\Big)+\theta^-\Big(\frac{\partial\phi'^n}{\partial\phi^m}\,\psi^m_-\Big)+\theta^+\theta^-\Big(\frac{\partial\phi'^n}{\partial\phi^m}\,F^m-\frac{\partial^2\phi'^n}{\partial\phi^m\partial\phi^l}\,\psi^m_+\psi^l_-\Big)\\
&:=\phi'^n+\theta^+\psi'^n_++\theta^-\psi'^n_-+\theta^+\theta^-F'^n
\end{aligned}
\]
特别地，费米子 $\psi^n_\pm$ 按全纯切向量变换。玻色子 $F^n$ 在靶空间上不按张量变换，但移位后的场 $F^n-\Gamma^n{}_{ml}\psi^m_+\psi^l_-$ 按全纯切向量变换。}我们先专注于手征多重态的场，
% [跨页接续: 原文第15页/第16页]
% ==== 原文 PDF 第 16 页 ====
反手征多重态的场可由共轭得到。首先，费米子 $\psi^n_\pm$ 自然被等同于 $\Sigma$ 上取值于全纯切丛拉回的左/右手旋量：$\psi_+\in\phi^*T^{(1,0)}M\otimes K_\Sigma^{1/2}$，$\psi_-\in\phi^*T^{(1,0)}M\otimes K_\Sigma^{-1/2}$。复玻色子 $F^n$ 在靶空间的坐标变换下不按张量变换，但 $F^n-\Gamma^n{}_{ml}\psi^m_+\psi^l_-$ 自然是全纯切丛拉回 $\phi^*T^{(1,0)}M$ 的一个截面。
% [原文疑误: 原文作 "under coordinate transformations of on the target"，多出一个介词 "of"，中译按「靶空间上的坐标变换」处理。]

我们要考虑的第二种超对称项只对费米子坐标的一半作积分，但超对称性限制了允许的被积式。特别地，若 $W(\Phi^n)$ 是手征超场 $\Phi^n$ 的全纯函数，则
\begin{equation}
\int d^2z\,d^2\theta\,W(\Phi^n)\big|_{\bar\theta^\pm=0}\,,
\tag{2.25}
\end{equation}
其中 $\int d^2\theta\,W|_{\bar\theta^\pm=0}$ 取出 $W$ 的 $\theta^+\theta^-$ 项；此式超对称不变，论证与 $D$ 项类似：对 $Q_\pm$ 的变分消去一个 $\theta^\pm$，故结果在费米子积分后为零；对 $\bar Q_\pm$ 的变分在费米子积分后不为零，但结果是全导数。这类表达式称为 $F$ 项，$W$ 称为超势。需要强调，$F$ 项关于 $\bar Q_\pm$ 的不变性依赖于一个事实：超势 $W$ 本身是手征超场，即手征超场的全纯函数。显式完成费米子积分后，超势/$F$ 项贡献
\begin{equation}
\int d^2z\,d^2\theta\,W(\Phi^n)\big|_{\bar\theta^\pm=0}=\int d^2z\big[F^nW_n-W_{nm}\psi^n_+\psi^m_-+\bar F^{\bar n}\bar W_{\bar n}-\bar W_{\bar n\bar m}\psi^n_-\psi^m_+\big]\,,
\tag{2.26}
\end{equation}
其中 $W_n=\partial_{\phi^n}W$，$W_{nm}=\partial_{\phi^n}\partial_{\phi^m}W$，依此类推。
% [Tutor Note: (2.26) 是显式算出 $\int d^2\theta$ 的结果：把 $\Phi^n=\phi^n+\theta^+\psi^n_++\theta^-\psi^n_-+\theta^+\theta^-F^n$ 代入 $W(\Phi^n)$ 并保留 $\theta^+\theta^-$ 分量，即得 $F^nW_n$ 与 Yukawa 项 $W_{nm}\psi^n_+\psi^m_-$；负号来自 $\psi$ 的反对易性与 $\theta$ 积分的定向约定。$\bar F^{\bar n}\bar W_{\bar n}$ 一项则是反手征部分的对应物。]
% [存疑: (2.26) 第二个费米子双线性写作 $\bar W_{\bar n\bar m}\psi^n_-\psi^m_+$，其费米子上指标 n、m 未加横线；若按"反手征部分由共轭得到"的惯例，此处似应出现 $\bar\psi$。原文可能只是省略了复共轭的记号（把 $\bar\psi^{\bar n}_-\bar\psi^{\bar m}_+$ 简写为 $\psi^n_-\psi^m_+$ 以强调组合方式），也可能漏加横线。此处照抄原样，供审校核对。]

类似地，可以构造超对称的扭 $F$ 项，即扭手征超场 $\tilde\Phi^{\tilde n}$ 的全纯函数 $\tilde W(\tilde\Phi^{\tilde n})$（扭超势）的积分：
\begin{equation}
\int d^2z\,d^2\tilde\theta\,\tilde W(\tilde\Phi^{\tilde n})\big|_{\bar\theta^+=\theta^-=0}\,,
\tag{2.27}
\end{equation}
其中 $\int d^2\tilde\theta\,\tilde W|_{\bar\theta^+=\theta^-=0}$ 取出 $\tilde W(\tilde\Phi^{\tilde n})$ 的 $\theta^+\bar\theta^-$ 项。

把 $D$ 项与 $F$ 项组合进手征（及反手征）多重态的理论，我们得到一族由凯勒流形 $M$ 与全纯超势 $W:M\to\mathbb{C}$ 标记的二维 $\mathcal{N}=(2,2)$ 理论：
\begin{equation}
S=\int d^2z\,d^4\theta\,K+\frac12\Big(\int d^2z\,d^2\theta\,W\big|_{\bar\theta^\pm=0}+\text{c.c.}\Big)\,.
\tag{2.28}
\end{equation}
% [Tutor Note: (2.28) 即二维 $\mathcal{N}=(2,2)$ 手征多重态的一般作用量：$D$ 项（凯勒势 $K$）给出靶空间 $M$ 上的非线性 $\sigma$ 模型，$F$ 项给出全纯超势 $W$。$F$ 项是全纯数据，受非重整化定理保护，故重整化群流只改变 $K$ 而不改变 $W$；这正是把 Landau--Ginzburg 模型与 Calabi--Yau 非线性 $\sigma$ 模型联系起来的 CY/LG 对应的技术基础。]
这些理论常被称为 Landau--Ginzburg 模型，但它们只是四维 Wess--Zumino 模型 [30, 31] 的维数约化。值得注意的是，复玻色子 $\bar F^{\bar n}$ 是辅助场，在用凯勒度规对偶化之后，其运动方程把它们限制为
\begin{equation}
\bar F_n-\Gamma_{n\bar m\bar l}\bar\psi^{\bar m}_-\bar\psi^{\bar l}_+=W_n\,.
\tag{2.29}
\end{equation}
写成（拉回的）全纯 1-形式，这就是 $\bar F=\partial_MW$，其中 $\partial_M=d\phi^n\partial_{\phi^n}$ 是 $M$ 上的全纯外微分。共轭场 $F^n$ 同样被限制：其运动方程（经移位与对偶化后）为 $F=\bar\partial_M\bar W$，其中 $\bar\partial_M$ 是 $M$ 上的反全纯外微分，即 Dolbeault 微分。
% [Tutor Note: 易错预警：(2.29) 左边写作 $\bar F_n$（下标），而正文讨论的是复玻色子 $\bar F^{\bar n}$（上标）；此外 $W_n=\partial_{\phi^n}W$ 是 $(1,0)$-形式的分量。原文所谓 "dualizing with the Kähler metric"（用凯勒度规对偶化）既包含对辅助场的对偶化（参见第 19 页脚注 19 对同一记号滥用的说明），也包含指标的升降。核对时以 (2.29) 的下标形式为准：它是 $1$-形式等式 $\bar F=\partial_MW$ 的分量表述。]

% ==== 原文 PDF 第 17 页 ====
可以追问：这个 Landau--Ginzburg 模型在什么条件下保持二维 $\mathcal{N}=(2,2)$ 代数的两个 $U(1)$ $R$ 对称性？$D$ 项与 $F$ 项必须各自独立地保持这些对称性，因为它们在矢量旋转或轴旋转下互不混合。就 $D$ 项而言，$d^4\theta$ 在轴与矢量 $R$ 对称旋转下都不变，因此只要能为手征超场指定一组荷，使 $K(\Phi^n,\bar\Phi^{\bar n})$ 在每个对称性下的总荷都为零，$D$ 项就同时保持两者。例如，若 $K$ 只是组合 $\Phi^n\bar\Phi^{\bar n}$ 的函数，它对任意荷分配都不变。

超势项更有意思。这里 $d^2\theta$ 在矢量 $R$ 对称性 $U(1)_V$ 下的荷为 $-2$，在轴 $U(1)$ 下的荷为 $0$。因此，为了保持矢量 $R$ 对称性，超势在矢量 $R$ 对称性 $U(1)_V$ 下的总荷必须为 $2$，而在轴 $R$ 对称性 $U(1)_A$ 下的荷为 $0$。就轴对称性而言，常见做法是把所有手征场都取 $U(1)_A$ 荷为零。就矢量对称性而言，若把超势取为拟齐次（次数为 $2$），则总荷为 $2$ 是可能的：存在 $U(1)_V$ 荷 $q_n$ 的某种取法，使得 $W(\lambda^{q_n}\Phi^n)=\lambda^2W(\Phi^n)$。

至此的分析都是经典的，但存在有趣的反常，可能阻止对称性在量子层面被保持，即路径积分测度可能并非不变。我们只陈述这项分析的结论，细节可参看（例如）[9]。由于费米子在 $R$ 对称性下带荷，必须研究路径积分测度中的费米子零模。结果是：超势不受微扰重整化，故对称性破坏的修正不可能在微扰论中产生；因此，只要选取（次数为 $2$ 的）拟齐次超势，矢量 $R$ 对称性 $U(1)_V$ 就是量子对称性。

轴 $R$ 对称性就没有这般幸运，它常常带有反常。应用漂亮的 Atiyah--Singer 指标定理，可以把结论几何地陈述为：凯勒流形 $M$ 必须有消失的第一陈类 $c_1(M)=0$。特别地，若 $\mathcal{X}$ 是 Calabi--Yau 流形、$W:\mathcal{X}\to\mathbb{C}$ 是拟齐次超势，则矢量与轴 $R$ 对称性都是无反常的对称性。反之，对具有一般凯勒靶空间与一般超势的 Landau--Ginzburg 模型，矢量/轴 $R$ 对称性会被破坏/成为反常的。
% [Tutor Note: 轴 $R$ 对称性的反常可由 Atiyah--Singer 指标定理读出：使路径积分测度在 $U(1)_A$ 下不变的条件正比于靶流形的第一陈类 $c_1(M)$，故 $c_1(M)=0$（Calabi--Yau 条件）是轴 $R$ 对称性在量子层面存活的前提；而矢量 $R$ 对称性因超势不受微扰重整化（全纯性 + 非重整化定理）而安全。这条"几何判据"是二维 $\mathcal{N}=(2,2)$ 理论中把 $R$ 反常与陈类联系起来的标准结果。]

本节最后指出，Landau--Ginzburg 理论的镜像对称不保持超势的存在性，也不保持非平凡曲率。例如，经典工作 [32] 的 2.2.2 节给出各种可能性的三个简单例子：以 $\mathbb{P}^{N-1}$ 为靶空间（无超势）的非线性 $\sigma$ 模型，其镜像是 $N-1$ 个手征场 $\Phi^n$（取值于 $\mathbb{C}^*$）、超势为 $W=\Phi^1+\dots+\Phi^{N-1}+q(\Phi^1\cdots\Phi^{N-1})^{-1}$ 的理论，也称 $A_{N-1}$ Toda 理论。\footnote{我们指出，这一例镜像对称在 [32] 之前已由拓扑扭变理论所知，参见例如 [33--38]。[32] 的工作把它推广到完整的物理理论：在 $\mathbb{P}^{N-1}$ 理论（的某个规范线性 $\sigma$ 模型实现）中求和瞬子贡献，从而找出其镜像 Landau--Ginzburg 模型。}后一例子尤其有趣，因为 $c_1(\mathbb{P}^{N-1})=N\neq0$，故该 $\sigma$ 模型具有反常的轴 $R$ 对称性 $U(1)_A$；其镜像对应 $A_{N-1}$ Toda 理论超势的非齐次性，因此该 Landau--Ginzburg 模型破坏 $U(1)_V$。\footnote{结果表明，这些 $R$ 对称性群中存在一个非平凡的离散子群在量子理论中保留下来。在 Landau--Ginzburg 模型中，$U(1)_V$ 的 $\mathbb{Z}_{2N}$ 子群不被破坏：在 $\Phi^n\to\zeta\Phi^n$ 下（$\zeta$ 为 $N$ 次单位根），超势齐次地（权为 $2$）变换，即 $\Phi^n$ 的 $\mathbb{Z}_{2N}$ 荷为 $2$。虽然不能用这个离散矢量 $R$ 对称性对 Landau--Ginzburg 模型作 $A$ 型扭变，但它确实把 $B$ 型扭变的 $\mathbb{Z}_2$ 分次精细化为 $\mathbb{Z}_{2N}$ 分次。这是 $\mathbb{P}^{N-1}$ $\sigma$ 模型的无反常 $\mathbb{Z}_{2N}$ 轴对称性、以及其 $A$ 型扭变上相应 $\mathbb{Z}_{2N}$ 分次的镜像。}
% [原文疑误: 脚注 17 中 "it does the refine the Z2 grading" 多出定冠词 "the"（且 "does refine" 才是强调用法）；另 "quasihomogenous" 一词漏字母 "e"，应为 "quasihomogeneous"。中译按正确读法处理。]
% [跨页接续: 脚注 17 起于原文第 17 页页脚，续于原文第 18 页页脚；此处按语义完整收录。]
% [Tutor Note: $A_{N-1}$ Toda 一例给出镜像对称中「超势的非齐次性 $\leftrightarrow$ 轴 $R$ 反常」的对应：$W=\Phi^1+\cdots+\Phi^{N-1}+q(\Phi^1\cdots\Phi^{N-1})^{-1}$ 只在 $\Phi^n\to\zeta\Phi^n$（$\zeta$ 为 $N$ 次单位根）下齐次，故 $\mathbb{P}^{N-1}$ 的 $\sigma$ 模型有反常的 $U(1)_A$，而其镜像 LG 模型的 $U(1)_V$ 被破坏。这是 CY/LG 对应中离散对称性与反常匹配的具体例证，也说明镜像对称可以把"几何上的非平凡曲率（$c_1\neq0$）"换成"超势的非齐次性"。]
